% ============================================================
%  PART II: APPENDICES
% ============================================================
\part{ضمائم و گسترش}
\label{part:appendices}

% ============================================================
%  APPENDIX 0: PRE-ASSEMBLY GOVERNANCE
% ============================================================
\chapter{پیوست نخست: مدیریت خلأ پیش از تشکیل مجلس مهستان}
\label{app:preassembly}

\section*{چرا این پیوست؟}

یکی از خطرناک‌ترین لحظات هر گذار دموکراتیک، \textbf{فاصله‌ی میان سقوط نظام پیشین
و تشکیل اولین نهاد منتخب} است. این فاصله --- که می‌تواند ۴ تا ۸ هفته به طول
بینجامد --- در صورت مدیریت نادرست می‌تواند به خلأ قدرت، هرج‌ومرج، یا مصادره‌ی
فرایند توسط یک جریان خاص منجر شود.

تجربه‌ی مصر (۲۰۱۱)، لیبی (۲۰۱۱) و یمن نشان می‌دهد که نبود پاسخ آماده برای این
فاصله، هزینه‌ای گاه غیرقابل‌جبران دارد.

\section*{مدل پیشنهادی: شورای هماهنگی گذار}

\begin{dangerbox}
\textbf{تعریف:} یک نهاد \textbf{موقت، کوچک، و با مأموریت محدود} که تنها برای مدیریت
این فاصله تشکیل می‌شود و پس از تشکیل مجلس مهستان، \textbf{بلافاصله منحل} می‌گردد.
\end{dangerbox}

\section*{اصول تشکیل شورای هماهنگی گذار}

\begin{bandbox}{اصل اول --- اتحاد و هماهنگی شرط ورود است}
شناخته‌شده‌ترین چهره‌های اپوزیسیون تنها در صورتی می‌توانند این نقش را بر عهده
بگیرند که \textbf{با یکدیگر به توافق رسیده باشند.} جریانی که از همکاری سرباز زند
یا ادعای انحصار داشته باشد، از ورود به این شورا محروم خواهد بود.
\end{bandbox}

\begin{bandbox}{اصل دوم --- نمایندگی متنوع}
ترکیب شورا باید طیف وسیعی از جریان‌های سیاسی، قومی، و جغرافیایی را دربر بگیرد.
حضور نمایندگان زنان، اقلیت‌های قومی، و جوانان الزامی است.
\end{bandbox}

\begin{bandbox}{اصل سوم --- شفافیت کامل}
تمام جلسات و تصمیمات شورا علنی است. هیچ تصمیمی با اکثریت کمتر از
\fraction{2}{3} اعضا اتخاذ نمی‌شود.
\end{bandbox}

% ★ FIX: single scaling method only — let TikZ handle it via `scale`
%   Removed the outer \resizebox wrapper that caused double-scaling
\begin{center}
\begin{tikzpicture}[
  scale=0.88, transform shape,
  every node/.style={align=right, font=\small},
  node distance=1cm                      % ★ tightened from 1.2cm
]

% ── Node 1: Vacuum ──
\node[
  rectangle, rounded corners=8pt,
  draw=MahestanRed, fill=DangerBg,
  text width=6cm, inner sep=8pt          % ★ unified width, cleaner padding
] (vacuum) {
  \textbf{\rl{خلأ پیش از مجلس مهستان}}\\[3pt]
  \rl{فاصله‌ی بین سقوط نظام}\\
  \rl{تا تشکیل مجلس مهستان}\\
  \rl{(تخمین: ۴ تا ۸ هفته)}
};

% ── Node 2: Council ──
\node[
  rectangle, rounded corners=8pt,
  draw=MahestanGold, fill=WarningBg,
  text width=6cm, inner sep=8pt,
  below=of vacuum
] (council) {
  \textbf{\rl{شورای هماهنگی گذار}}\\[4pt]
  \rl{رهبران شناخته‌شده‌ی اپوزیسیون}\\
  \rl{به شرط هماهنگی و اتحاد}\\[3pt]
  {\footnotesize
    \rl{حداکثر ۱۵–۲۰ نفر}\\
    \rl{نمایندگی جریان‌های اصلی}%
  }
};

% ── Node 3: Mandate ──
\node[
  rectangle, rounded corners=6pt,
  draw=MahestanBlue, fill=NodeBg,
  text width=6cm, inner sep=8pt,         % ★ was 7cm — now consistent
  below=of council
] (mandate) {
  \textbf{\rl{مأموریت محدود}}\\[3pt]
  \rl{\textcolor{MahestanGreen}{\checkmark{} مجاز:}}\\
  \rl{برگزاری انتخابات مجلس مهستان}\\
  \rl{صدور ۱۰ فرمان اضطراری}\\
  \rl{تثبیت نظم عمومی}\\
  \rl{اطلاع‌رسانی به جامعه‌ی بین‌المللی}\\[3pt]
  \rl{\textcolor{MahestanRed}{\crossmark{} غیرمجاز:}}\\
  \rl{تصمیمات ساختاری یا قانون‌گذاری}\\
  \rl{معاهدات بین‌المللی الزام‌آور}
};

% ── Node 4: Assembly ──
\node[
  rectangle, rounded corners=6pt,
  draw=MahestanGreen, fill=SuccessBg,
  text width=6cm, inner sep=8pt,
  below=of mandate
] (assembly) {
  \textbf{\rl{تشکیل مجلس مهستان}}
};

% ── Node 5: Dissolve ──
\node[
  rectangle, rounded corners=6pt,
  draw=MahestanGreen!70!black, fill=SuccessBg,
  text width=6cm, inner sep=6pt,
  below=of assembly
] (dissolve) {
  \textbf{\rl{انحلال فوری شورای هماهنگی گذار}}
};

% ── Arrows ──
\draw[maharrow=MahestanRed]   (vacuum)  -- (council);
\draw[maharrow=MahestanGold]  (council) -- (mandate);
\draw[maharrow=MahestanBlue]  (mandate) -- (assembly);
\draw[maharrow=MahestanGreen] (assembly) -- (dissolve);

\end{tikzpicture}
\end{center}



\begin{bandbox}{اصل چهارم --- محدودیت مأموریت}
شورا هیچ صلاحیتی فراتر از «اداره‌ی ضروریات» و «برگزاری انتخابات» ندارد.
\end{bandbox}

\begin{bandbox}{اصل پنجم --- نظارت بین‌المللی}
از همان روز اول، ناظران بین‌المللی (سازمان ملل، اتحادیه اروپا، \lr{OSCE}) دعوت
می‌شوند تا بر کار شورا نظارت کنند.
\end{bandbox}

\begin{notebox}[title={\rl{فراخوان به جریان‌های اپوزیسیون}}]
آمادگی برای این فاز باید از هم‌اکنون آغاز شود. هر جریانی که می‌خواهد نقشی در
مدیریت این مرحله داشته باشد، باید:
\begin{enumerate}
  \item نماینده‌ی مشخص برای مذاکره با سایر جریان‌ها معرفی کند
  \item خطوط قرمز و نقاط توافق خود را صریحاً اعلام کند
  \item پیش‌نویس توافق‌نامه‌ی همکاری را آماده کند
  \item درباره‌ی ترکیب شورا و آیین‌نامه‌ی تصمیم‌گیری از هم‌اکنون مذاکره کند
\end{enumerate}
\end{notebox}

% ============================================================
%  APPENDIX A: COMPARATIVE ANALYSIS
% ============================================================
\chapter{ضمیمه‌ی الف: تحلیل تطبیقی گذارهای دموکراتیک}
\label{app:comparative}

\section{الگوهای موفق}

\vspace{0.5\baselineskip}
\noindent
\begin{adjustbox}{max width=\linewidth}
\begin{tabular}{@{}p{2.5cm}p{3cm}p{4cm}p{4cm}@{}}
\toprule
\textbf{کشور} & \textbf{دوره} & \textbf{عوامل موفقیت} & \textbf{درس برای ایران} \\
\midrule
\rowcolor{SuccessBg}
آفریقای جنوبی & ۱۹۹۴--۱۹۹۶ &
  مذاکرات \lr{CODESA}؛ دو قانون اساسی؛ مشاوره‌ی ۲ میلیون نفری &
  مدل دو قانون اساسی قابل اجرا در ایران است \\
\addlinespace
اسپانیا & ۱۹۷۵--۱۹۷۸ &
  پیمان مونکلوا؛ اجماع نخبگان؛ تصویب قانون اساسی با ۸۸\% آرا &
  اجماع نخبگان حتی با رقیبان ایدئولوژیک \\
\addlinespace
\rowcolor{SuccessBg}
لهستان & ۱۹۸۹ &
  مذاکرات میزگرد؛ انتخابات تدریجی؛ حفظ نهادها &
  «اصلاح نه تخریب» جواب داد \\
\addlinespace
شیلی & ۱۹۸۸--۱۹۹۰ &
  رفراندوم مسالمت‌آمیز؛ حفظ ارتش در چارچوب &
  عبور از دیکتاتور بدون جنگ داخلی \\
\bottomrule
\end{tabular}
\end{adjustbox}

\section{ضد الگوها (از چه باید پرهیز کرد)}

\vspace{0.5\baselineskip}
\noindent
\begin{adjustbox}{max width=\linewidth}
\begin{tabular}{@{}p{2.5cm}p{3cm}p{4cm}p{4cm}@{}}
\toprule
\textbf{کشور} & \textbf{اشتباه کلیدی} & \textbf{پیامد} & \textbf{هشدار برای ایران} \\
\midrule
\rowcolor{DangerBg}
عراق & انحلال کامل ارتش و حزب بعث &
  فروپاشی نهادها؛ شکل‌گیری \lr{ISIS} &
  حفظ ساختارهای غیرایدئولوژیک ضروری است \\
\addlinespace
\rowcolor{DangerBg}
لیبی & نبود توافق پیش از سقوط قذافی &
  جنگ داخلی؛ تجزیه &
  برنامه‌ریزی پیشین برای مرحله‌ی گذار \\
\addlinespace
\rowcolor{DangerBg}
مصر & ارتش جای سیاست‌مداران را گرفت &
  بازگشت اقتدارگرایی با ظاهر دموکراتیک &
  نظامی‌ها نباید رهبری سیاسی گذار را بگیرند \\
\bottomrule
\end{tabular}
\end{adjustbox}

% ── Detailed Case Studies ──────────────────────────────────
\section{تحلیل تفصیلی: موردهای کلیدی}

\subsection{الف-۱. تونس (۲۰۱۱–۲۰۱۴): نزدیک‌ترین الگو}

تونس موفق‌ترین نمونه‌ی گذار دموکراتیک در خاورمیانه و شمال آفریقا محسوب می‌شود.

\needspace{4\baselineskip}
\begin{bandbox}{شباهت‌ها با طرح پیشنهادی}
\begin{itemize}
  \item انتخاب \textbf{مجلس مؤسسان ملی} (\lr{Assemblée nationale constituante}) در اکتبر ۲۰۱۱ به‌عنوان نهاد محوری گذار
  \item این مجلس هم وظیفه‌ی \textbf{قانون‌گذاری موقت} و هم \textbf{تدوین قانون اساسی} را بر عهده داشت
  \item تشکیل \textbf{دولت تکنوکرات} تحت نظارت مجلس
  \item استفاده از \textbf{ائتلاف سه‌جانبه} (ترویکا) برای مدیریت تکثرگرایانه
\end{itemize}
\end{bandbox}

\begin{notebox}[title={\rl{دروس آموخته}}]
\begin{itemize}
  \item \textbf{بحران ۲۰۱۳}: «چهارگانه‌ی گفتگوی ملی» (شامل اتحادیه‌ی کارگری \lr{UGTT}، اتحادیه‌ی کارفرمایان، کانون وکلا و جامعه حقوق بشر) از سقوط فرایند جلوگیری کرد. \textbf{درس: نهادهای میانجی‌گر غیردولتی حیاتی‌اند.}
  \item تأخیر در تدوین قانون اساسی (سه سال به جای یک سال): \textbf{واقع‌بینی در زمان‌بندی ضروری است.}
  \item عدم رسیدگی کافی به اصلاحات اقتصادی و بخش امنیتی، زمینه‌ی بازگشت استبداد را فراهم کرد (۲۰۲۱).
\end{itemize}

\textbf{پیشنهاد}: مجلس مهستان باید از همان ابتدا \textbf{کمیسیون مستقل اصلاحات اقتصادی} و \textbf{کمیسیون اصلاح بخش امنیتی} را تشکیل دهد. نهادهای مستقل جامعه‌ی مدنی باید نقش رسمی مشورتی و نظارتی داشته باشند.
\end{notebox}

\subsection{الف-۲. آفریقای جنوبی (۱۹۹۰–۱۹۹۶): مدل عدالت انتقالی}

\begin{bandbox}{فرایند گذار}
\begin{itemize}
  \item \textbf{مذاکرات کادسا} (\lr{CODESA I \& II}) --- ۱۹۹۱–۱۹۹۲
  \item \textbf{مجمع مذاکره‌ی چندحزبی} (\lr{MPNF}) --- ۱۹۹۳
  \item \textbf{شورای اجرایی انتقالی} (\lr{TEC}) --- نظارت بر انتخابات ۱۹۹۴
  \item \textbf{قانون اساسی موقت} $\leftarrow$ \textbf{انتخابات} $\leftarrow$ \textbf{مجلس مؤسسان} $\leftarrow$ \textbf{قانون اساسی نهایی} (۱۹۹۶)
  \item \textbf{کمیسیون حقیقت و آشتی} (\lr{TRC}) به ریاست دزموند توتو
\end{itemize}
\end{bandbox}

\begin{notebox}[title={\rl{دروس آموخته}}]
\begin{itemize}
  \item مدل \textbf{دو مرحله‌ای قانون اساسی} (موقت + نهایی) بسیار مؤثر بود.
  \item کمیسیون حقیقت و آشتی با \textbf{عفو مشروط} (در ازای اعتراف کامل)، مسیر میانه‌ای بین انتقام و فراموشی ایجاد کرد.
  \item \textbf{ضمانت‌های حقوق اقلیت‌ها} در قانون اساسی موقت، ترس اقلیت سفیدپوست را کاهش داد.
\end{itemize}

\textbf{پیشنهاد}: یک \textbf{منشور اصول بنیادین غیرقابل‌نقض} باید پیش از انتخابات مجلس مهستان تدوین شود.
\end{notebox}

\subsection{الف-۳. اسپانیا (۱۹۷۵–۱۹۸۲): مدل گذار تدریجی}

\begin{bandbox}{ویژگی‌های کلیدی}
\begin{itemize}
  \item گذار \textbf{مذاکره‌شده و تدریجی} پس از مرگ فرانکو
  \item \textbf{پیمان‌های مونکلوآ} (۱۹۷۷) --- توافق اقتصادی‌اجتماعی بین احزاب، کارگران و کارفرمایان
  \item تصویب قانون اساسی ۱۹۷۸ در رفراندوم با ۸۸\% آرا
\end{itemize}
\end{bandbox}

\begin{notebox}[title={\rl{درس کلیدی}}]
\textbf{قانون فراموشی} (\lr{Pacto del Olvido}) اگرچه در کوتاه‌مدت ثبات آورد، در بلندمدت مشکل‌ساز شد. \textbf{درس: عدالت انتقالی را نمی‌توان به‌طور کامل نادیده گرفت.} کودتای نافرجام ۱۹۸۱ نشان داد که \textbf{اصلاح بخش نظامی} از همان ابتدا حیاتی است.
\end{notebox}

\subsection{الف-۴. لهستان (۱۹۸۹): مدل میز مذاکره}

\begin{bandbox}{ویژگی‌های کلیدی}
\begin{itemize}
  \item \textbf{مذاکرات میزگرد} بین جنبش همبستگی (\lr{Solidarność}) و حکومت کمونیست
  \item \textbf{انتخابات نیمه‌آزاد} ژوئن ۱۹۸۹ --- دولت ائتلافی به رهبری مازوویتسکی
  \item سیاست \textbf{«خط ضخیم»} مازوویتسکی: جدایی از گذشته بدون انتقام‌جویی
\end{itemize}

\textbf{پیشنهاد}: برنامه‌ی اقتصادی دوران گذار باید \textbf{تدریجی و با شبکه‌ی ایمنی اجتماعی} همراه باشد. شوک‌درمانی اقتصادی می‌تواند مشروعیت دموکراسی نوپا را تضعیف کند.
\end{bandbox}

\subsection{الف-۵. شیلی (۱۹۸۸–۱۹۹۰): مدل رفراندوم و ائتلاف}

\begin{bandbox}{ویژگی‌ها و درس کلیدی}
\begin{itemize}
  \item \textbf{رفراندوم ۱۹۸۸}: «نه» به پینوشه --- \textbf{ائتلاف احزاب} (\lr{Concertación}) --- ۱۷ حزب با ایدئولوژی‌های متفاوت
  \item توانایی جریان‌های مختلف برای کنار گذاشتن اختلافات و تشکیل ائتلاف فراگیر، عامل اصلی موفقیت بود
  \item فرایند عدالت انتقالی \textbf{تدریجی} بود --- کمیسیون حقیقت (گزارش رتیگ ۱۹۹۱) ابتدا فقط حقیقت‌یابی کرد
\end{itemize}
\end{bandbox}

\subsection{الف-۶. نمونه‌های ناموفق --- درس‌های منفی}

\vspace{0.5\baselineskip}
\begin{adjustbox}{max width=\linewidth}
\begin{tabular}{@{}p{2.5cm}p{4cm}p{6.5cm}@{}}
\toprule
\textbf{کشور} & \textbf{مشکل اصلی} & \textbf{درس برای ایران} \\
\midrule
\rowcolor{DangerBg}
لیبی (۲۰۱۱--) &
  عدم وجود نهاد محوری مورد توافق؛ شوراهای رقیب؛ مداخله‌ی خارجی &
  نهاد واحد گذار با مشروعیت انتخاباتی ضروری است \\
\addlinespace
\rowcolor{DangerBg}
مصر (۲۰۱۱--۲۰۱۳) &
  ارتش به‌عنوان «ضامن» گذار؛ تقابل اسلامگرا-سکولار &
  استقلال نهاد گذار از نیروهای نظامی و جلوگیری از دوقطبی‌سازی \\
\addlinespace
\rowcolor{DangerBg}
عراق (۲۰۰۳--) &
  اشغال خارجی؛ سیستم محاصصه‌ای؛ انحلال ارتش &
  عدم انحلال کامل نهادهای اداری‌نظامی؛ اصلاح به‌جای تخریب \\
\addlinespace
\rowcolor{DangerBg}
یمن (۲۰۱۱--) &
  گذار مصنوعی و ظاهری؛ مذاکره‌ی بدون مشارکت عمومی &
  مشارکت واقعی مردم از طریق انتخابات \\
\bottomrule
\end{tabular}
\end{adjustbox}

\subsection{الف-۷. جدول مقایسه‌ای جامع}

\vspace{0.5\baselineskip}
\begin{landscape}
\thispagestyle{plain}
\begin{adjustbox}{max width=\linewidth}
\begin{tabular}{@{}p{2.8cm}p{3.5cm}p{2.5cm}p{2.5cm}p{2.5cm}p{3.5cm}@{}}
\toprule
\textbf{معیار} & \textbf{تونس} & \textbf{آفریقای جنوبی} & \textbf{اسپانیا} & \textbf{لهستان} & \textbf{طرح مهستان} \\
\midrule
\rowcolor{MahestanLightBlue!30}
\textbf{نهاد محوری} &
  مجلس مؤسسان منتخب &
  شورای اجرایی انتقالی + مجلس مؤسسان &
  پارلمان منتخب &
  دولت ائتلافی &
  مجلس مهستان منتخب \\
\addlinespace
\textbf{مدت گذار} & ۳ سال & ۶ سال & ۳ سال & ۲ سال & حداکثر ۲ سال \\
\addlinespace
\rowcolor{MahestanLightBlue!30}
\textbf{سیستم انتخاباتی} &
  فهرستی نسبی &
  فهرستی نسبی &
  منطقه‌ای نسبی &
  مختلط &
  فهرستی ترکیبی \\
\addlinespace
\textbf{عدالت انتقالی} &
  \lr{IVD} (۲۰۱۴) &
  \lr{TRC} (۱۹۹۵) &
  قانون فراموشی (ناقص) &
  «خط ضخیم» &
  پیش‌بینی‌شده در طرح \\
\addlinespace
\rowcolor{MahestanLightBlue!30}
\textbf{نقش نظامیان} &
  حامی گذار &
  مذاکره‌شده &
  چالش‌برانگیز &
  زیر نظر دولت &
  تحت نظارت مجلس \\
\addlinespace
\textbf{قانون اساسی} &
  مجلس مؤسسان + رفراندوم &
  دو مرحله‌ای &
  رفراندوم &
  اجماع + رفراندوم &
  تدوین + رفراندوم \\
\addlinespace
\rowcolor{SuccessBg}
\textbf{نتیجه} &
  بازگشت استبداد (۲۰۲۱) &
  دموکراسی پایدار &
  دموکراسی پایدار &
  دموکراسی (با چالش) &
  هدف: دموکراسی پایدار \\
\bottomrule
\end{tabular}
\end{adjustbox}
\end{landscape}

% ============================================================
%  APPENDIX B: TIMELINE TABLE
% ============================================================
\chapter{ضمیمه‌ی ب: جدول زمانی پیشنهادی اجرا}
\label{app:timeline}

\begin{notebox}
تمام زمان‌بندی‌های زیر \textbf{نسبی} هستند. «هفته‌ی ۱» یعنی هفته‌ی اول پس از
\textbf{روز صفر} (آغاز رسمی گذار). مجلس مهستان می‌تواند این جدول را بر اساس
شرایط واقعی تعدیل کند.
\end{notebox}

\subsection*{فاز صفر: آمادگی (پیش از گذار)}
\begin{adjustbox}{max width=\linewidth}
\begin{tabular}{@{}lp{10cm}@{}}
\toprule
\textbf{زمان} & \textbf{اقدام} \\
\midrule
\rowcolor{MahestanLightGray}
مستمر & شکل‌گیری خودجوش جریان‌ها و فهرست‌ها \\
مستمر & تبلیغات، معرفی کاندیداها و ائتلاف‌سازی \\
\rowcolor{MahestanLightGray}
مستمر & طراحی زیرساخت‌های فنی رأی‌گیری (بستر دیجیتال) \\
مستمر & تدوین منشور اصول بنیادین غیرقابل‌نقض \\
\rowcolor{MahestanLightGray}
مستمر & تهیه‌ی فهرست قوانین فوری (تعلیق/الغای فوری) \\
\bottomrule
\end{tabular}
\end{adjustbox}

\subsection*{فاز یک: تثبیت اولیه (هفته‌ی ۱--۴)}
\begin{adjustbox}{max width=\linewidth}
\begin{tabular}{@{}lp{10cm}@{}}
\toprule
\textbf{هفته} & \textbf{اقدام} \\
\midrule
\rowcolor{WarningBg}
۱ & اعلام رسمی آغاز دوره‌ی گذار؛ انتصاب مدیران موقت فرمانداری‌ها \\
۱--۲ & آغاز ثبت‌نام رسمی نامزدها و فهرست‌ها \\
\rowcolor{WarningBg}
۱--۲ & تشکیل ستاد مرکزی انتخابات با مشارکت ناظران بین‌المللی \\
۲--۳ & دوره‌ی تبلیغات رسمی \\
\rowcolor{WarningBg}
۳--۴ & دعوت از ناظران بین‌المللی (سازمان ملل، اتحادیه اروپا) \\
\bottomrule
\end{tabular}
\end{adjustbox}

\subsection*{فاز دو: انتخابات و تشکیل مجلس (هفته‌ی ۴--۸)}
\begin{adjustbox}{max width=\linewidth}
\begin{tabular}{@{}lp{10cm}@{}}
\toprule
\textbf{هفته} & \textbf{اقدام} \\
\midrule
\rowcolor{DangerBg}
۴--۵ & برگزاری انتخابات مجلس مهستان \\
۵--۶ & شمارش آرا و تأیید نتایج \\
\rowcolor{DangerBg}
۶ & تشکیل رسمی مجلس مهستان \\
۶--۷ & انتخاب هیأت رئیسه و تعیین کمیسیون‌ها \\
\rowcolor{DangerBg}
۷--۸ & تشکیل کابینه‌ی اجرایی و رأی اعتماد \\
\bottomrule
\end{tabular}
\end{adjustbox}

% ============================================================
%  APPENDIX C: RISK ANALYSIS
% ============================================================
\chapter{ضمیمه‌ی ج: تحلیل ریسک و راهکارهای مقابله}
\label{app:risk}

\section{جدول ریسک‌های اصلی}

\vspace{0.5\baselineskip}
\begin{adjustbox}{max width=\linewidth}
\begin{tabular}{@{}p{0.4cm}p{3cm}p{1.8cm}p{1.5cm}p{6cm}@{}}
\toprule
\textbf{\#} & \textbf{ریسک} & \textbf{احتمال} & \textbf{شدت} & \textbf{راهکار پیشنهادی} \\
\midrule
\rowcolor{DangerBg}
۱ & خلأ امنیتی پس از سقوط نظام &
  بالا & بسیار بالا &
  تشکیل فوری شورای امنیتی موقت؛ حفظ ساختار پایه‌ای نیروهای انتظامی با رهبری جدید؛ کمک‌خواهی از نهادهای بین‌المللی \\
\addlinespace
\rowcolor{DangerBg}
۲ & تجزیه‌طلبی و درگیری‌های قومی &
  متوسط & بسیار بالا &
  تضمین حقوق قومی و زبانی در منشور؛ کرسی‌های تضمینی در مجلس؛ فدرالیسم به‌عنوان گزینه \\
\addlinespace
\rowcolor{DangerBg}
۳ & بازگشت اقتدارگرایی (کودتا) &
  متوسط & بسیار بالا &
  نظارت مجلس بر نیروهای نظامی؛ پراکنده‌سازی قدرت نظامی؛ حمایت بین‌المللی \\
\addlinespace
\rowcolor{WarningBg}
۴ & بحران اقتصادی حاد &
  بالا & بالا &
  تشکیل فوری تیم مدیریت بحران؛ درخواست کمک از \lr{IMF}/\lr{WB}؛ تثبیت ارزی \\
\addlinespace
\rowcolor{WarningBg}
۵ & مداخله‌ی خارجی &
  متوسط & بالا &
  دیپلماسی فعال و متنوع؛ عدم وابستگی به یک قدرت؛ شفافیت در روابط خارجی \\
\addlinespace
\rowcolor{WarningBg}
۶ & اختلاف درون مجلس و بن‌بست سیاسی &
  بالا & متوسط &
  سازوکارهای حل اختلاف (میانجی‌گری)؛ نقش ناظران جامعه‌ی مدنی \\
\addlinespace
\rowcolor{NodeBg}
۷ & مصادره‌ی جنبش توسط یک جریان &
  متوسط & بالا &
  سیستم فهرستی نسبی؛ آستانه‌ی بالای تغییر قوانین (۷۵\%)؛ منشور غیرقابل‌نقض \\
\addlinespace
\rowcolor{NodeBg}
۸ & خشونت انتقام‌جویانه &
  بالا & بالا &
  قانون عدالت انتقالی فوری؛ محاکم قانونی؛ کمپین‌های ملی آشتی \\
\addlinespace
\rowcolor{NodeBg}
۹ & تقلب در انتخابات &
  پایین (با نظارت) & متوسط &
  ناظران بین‌المللی؛ بستر دیجیتال شفاف \\
\addlinespace
\rowcolor{NodeBg}
۱۰ & فساد در نهادهای جدید &
  متوسط & متوسط &
  واحد بازرسی مستقل؛ شفافیت مالی نمایندگان؛ نظارت رسانه‌ای \\
\bottomrule
\end{tabular}
\end{adjustbox}

\section{سناریوهای بحرانی و پاسخ‌ها}

\begin{bandbox}{سناریو ۱ --- عدم تشکیل ائتلاف در مجلس}
\begin{itemize}
  \item فعال‌سازی سازوکار بند ۱۱ (رأی‌گیری فردی برای هر پست)
  \item میانجی‌گری توسط شخصیت‌های مورد احترام ملی
  \item در صورت عدم موفقیت: رأی‌گیری مجدد ظرف ۶۰ روز
\end{itemize}
\end{bandbox}

\begin{bandbox}{سناریو ۲ --- تلاش نظامی برای کودتا}
\begin{itemize}
  \item فراخوان فوری مقاومت مدنی از سوی مجلس
  \item فعال‌سازی پروتکل‌های بین‌المللی (بیانیه‌ی محکومیت، تحریم)
  \item جانشین فرمانده به‌عنوان فرمانده‌ی قانونی شناخته می‌شود
\end{itemize}
\end{bandbox}

\begin{bandbox}{سناریو ۳ --- بحران اقتصادی حاد}
\begin{itemize}
  \item تشکیل کمیته‌ی بحران اقتصادی با اختیارات ویژه
  \item درخواست کمک اضطراری از صندوق بین‌المللی پول و بانک جهانی
  \item اجرای برنامه‌ی یارانه‌ی مستقیم به شهروندان
\end{itemize}
\end{bandbox}

% ============================================================
%  APPENDIX D: ECONOMIC GOVERNANCE
% ============================================================
\chapter{ضمیمه‌ی د: حکمرانی اقتصادی در دوران گذار}
\label{app:economic}

\section{اصول راهنما}

\begin{notebox}
تجربه‌ی جهانی نشان می‌دهد که \textbf{بحران اقتصادی} بزرگ‌ترین تهدید برای دموکراسی‌های نوپاست. مردمی که گرسنه‌اند و کاری ندارند، به دموکراسی بی‌اعتماد می‌شوند. دموکراسی‌های موفق (اسپانیا، لهستان، کره‌ی جنوبی) همه از ثبات اقتصادی نسبی در دوران گذار برخوردار بودند.
\end{notebox}

\begin{enumerate}
  \item \textbf{تثبیت اقتصادی} مقدم بر اصلاحات ساختاری است
  \item \textbf{شبکه‌ی ایمنی اجتماعی} باید از همان روز اول فعال باشد
  \item \textbf{شفافیت} در دارایی‌های عمومی و بودجه‌ی دولت ضروری است
  \item \textbf{استفاده از تخصص} اقتصاددانان ایرانی داخل و خارج
\end{enumerate}

\section{اقدامات فوری (۱۰۰ روز اول)}

\begin{adjustbox}{max width=\linewidth}
\begin{tabular}{@{}lp{10.5cm}@{}}
\toprule
\textbf{اولویت} & \textbf{اقدام} \\
\midrule
\rowcolor{DangerBg}
فوری &
  حسابرسی فوری بانک مرکزی و صندوق توسعه‌ی ملی \\
\rowcolor{DangerBg}
فوری &
  تثبیت نرخ ارز و جلوگیری از فرار سرمایه \\
\rowcolor{DangerBg}
فوری &
  ادامه‌ی پرداخت حقوق کارمندان دولت و بازنشستگان \\
\rowcolor{WarningBg}
بالا &
  مذاکره با نهادهای مالی بین‌المللی برای بسته‌ی حمایتی \\
\rowcolor{WarningBg}
بالا &
  آزادسازی دارایی‌های مسدودشده (پس از رفع تحریم‌ها) \\
\rowcolor{NodeBg}
متوسط &
  رسیدگی به دارایی‌های غیرقانونی بنیادها و نهادهای حکومتی \\
\bottomrule
\end{tabular}
\end{adjustbox}

\section{اقدامات میان‌مدت (۱۰۰ روز تا ۲ سال)}

\begin{itemize}
  \item بازسازی شفاف بودجه‌ی ملی
  \item اصلاح تدریجی نظام یارانه‌ای با حمایت هدفمند از اقشار آسیب‌پذیر
  \item بازنگری قراردادهای نفتی و بین‌المللی
  \item آغاز خصوصی‌سازی شفاف نهادهای اقتصادی وابسته به حکومت (سپاه، بنیادها)
  \item ایجاد فضای جذب سرمایه‌گذاری خارجی
\end{itemize}

\begin{summarybox}
\textbf{درس از لهستان}: برنامه‌ی اقتصادی دوران گذار باید \textbf{تدریجی و با شبکه‌ی ایمنی اجتماعی} همراه باشد. شوک‌درمانی اقتصادی (الگوی لهستان) می‌تواند در بلندمدت مؤثر باشد اما در کوتاه‌مدت \textbf{درد اجتماعی شدیدی} ایجاد می‌کند که ممکن است مشروعیت دموکراسی نوپا را تضعیف کند. ایران باید مسیری میانه را بپیماید: اصلاحات اقتصادی ضروری، اما با سرعتی که مردم بتوانند تحمل کنند.
\end{summarybox}

% ============================================================
%  APPENDIX: TRANSITIONAL JUSTICE (condensed)
% ============================================================
\chapter{ضمیمه‌ی ه: چارچوب عدالت انتقالی}
\label{app:tj}

\section{مقایسه‌ی مدل‌های عدالت انتقالی}

\vspace{0.5\baselineskip}
\noindent
\begin{adjustbox}{max width=\linewidth}
\begin{tabular}{@{}p{3cm}lp{4.5cm}p{4.5cm}@{}}
\toprule
\textbf{کشور} & \textbf{مدل} & \textbf{نقاط قوت} & \textbf{نتیجه} \\
\midrule
\rowcolor{SuccessBg}
آفریقای جنوبی & \lr{TRC} (1995--2002) &
  عفو مشروط + حقیقت‌یابی &
  آشتی نسبی؛ دموکراسی پایدار \\
\addlinespace
شیلی & تدریجی &
  \lr{CONADEP}؛ غرامت؛ موزه‌ی حافظه &
  الگوی موفق محاکمات \\
\addlinespace
\rowcolor{SuccessBg}
آرژانتین & عدالت‌محور &
  \lr{CONADEP}؛ محاکمات گسترده &
  الگوی محاکمات جنایی \\
\addlinespace
\rowcolor{DangerBg}
اسپانیا & پیمان فراموشی &
  ثبات سیاسی کوتاه‌مدت &
  زخم‌های باز تاریخی \\
\addlinespace
\rowcolor{DangerBg}
تونس & \lr{IVD} (2014--2018) &
  کمیسیون حقیقت‌یابی &
  بازگشت استبداد (2021) \\
\bottomrule
\end{tabular}
\end{adjustbox}

\section{درس کلیدی}

\begin{summarybox}
\begin{center}
\begin{tikzpicture}
\node[
  rectangle, rounded corners=6pt,
  draw=MahestanRed, fill=DangerBg,
  text width=4.5cm, inner sep=8pt, font=\small
] (justice_only) {\textbf{\rl{فقط عدالت کیفری}}\\[4pt]\rl{خشم و انتقام باقی ماند}\\[2pt]\rl{\tiny (رواندا)}};

\node[
  rectangle, rounded corners=6pt,
  draw=MahestanGold, fill=WarningBg,
  text width=4.5cm, inner sep=8pt, font=\small,
  right=0.8cm of justice_only
] (peace_only) {\textbf{\rl{فقط آشتی}}\\[4pt]\rl{بی‌عدالتی نهادینه شد}\\[2pt]\rl{\tiny (اسپانیا)}};

\node[
  rectangle, rounded corners=6pt,
  draw=MahestanGreen, fill=SuccessBg,
  text width=4.5cm, inner sep=8pt, font=\small,
  right=0.8cm of peace_only
] (both) {\textbf{\rl{عدالت + آشتی}}\\[4pt]\rl{صلح پایدار}\\[2pt]\rl{\tiny (آفریقای جنوبی، شیلی)}};

\node[below=0.3cm of justice_only, font=\Large\color{MahestanRed}] {\lr{(x)}};
\node[below=0.3cm of peace_only, font=\Large\color{MahestanGold}] {\lr{(x)}};
\node[below=0.3cm of both, font=\Large\color{MahestanGreen}] {\lr{(v)}};
\end{tikzpicture}
\end{center}
\end{summarybox}

\section{مدل چهارستونی سازمان ملل}

\begin{notebox}[title={\rl{مدل کلاسیک \lr{UN Guidance Note 2010}}}]
مدل چهارستونی، استاندارد مرجع سازمان ملل برای عدالت انتقالی است و در اکثر گذارهای دموکراتیک دهه‌های اخیر به‌کار رفته است.
\end{notebox}

\vspace{0.5\baselineskip}
\begin{center}
\resizebox{0.92\linewidth}{!}{%
\begin{tikzpicture}[
  every node/.style={align=right, font=\small},
  node distance=0.8cm and 1.5cm
]
\node[
  rectangle, rounded corners=8pt,
  draw=MahestanPurple, fill=MahestanLightPurple,
  text width=5cm, inner sep=10pt, font=\bfseries\small
] (center) {
  \rl{عدالت انتقالی ایران}\\[4pt]
  \rl{\footnotesize \lr{Transitional Justice}}\\[2pt]
  \rl{\footnotesize مدل نهایی را مجلس مهستان}\\
  \rl{\footnotesize با مشارکت مردم تعیین می‌کند}
};

\node[mahbox=MahestanBlue, above right=1.5cm and 1.8cm of center] (p1) {
  \textbf{\rl{ستون ۱: حقیقت‌یابی}}\\[3pt]
  \rl{کمیسیون حقیقت و آشتی}\\
  \rl{مستندسازی جنایات}\\
  \rl{ثبت شهادت قربانیان}\\
  \rl{گزارش نهایی جامع}\\[2pt]
  \rl{\scriptsize الگو: \lr{TRC} آفریقای جنوبی}
};

\node[mahbox=MahestanRed, below right=1.5cm and 1.8cm of center] (p2) {
  \textbf{\rl{ستون ۲: عدالت کیفری}}\\[3pt]
  \rl{محاکمات عادلانه و علنی}\\
  \rl{تفکیک سطوح مسئولیت}\\
  \rl{دادرسی عادلانه}\\
  \rl{عفو مشروط (سطوح پایین)}\\[2pt]
  \rl{\scriptsize الگو: دادگاه‌های بین‌المللی}
};

\node[mahbox=MahestanOrange, above left=1.5cm and 1.8cm of center] (p3) {
  \textbf{\rl{ستون ۳: جبران خسارت}}\\[3pt]
  \rl{غرامت مالی فردی}\\
  \rl{بازگرداندن اموال مصادره‌شده}\\
  \rl{خدمات درمانی‌روانشناختی}\\
  \rl{اعاده‌ی حیثیت}\\[2pt]
  \rl{\scriptsize الگو: برنامه‌ی غرامت شیلی}
};

\node[mahbox=MahestanGreen, below left=1.5cm and 1.8cm of center] (p4) {
  \textbf{\rl{ستون ۴: تضمین عدم تکرار}}\\[3pt]
  \rl{اصلاح بخش امنیتی (\lr{SSR})}\\
  \rl{اصلاح دستگاه قضایی}\\
  \rl{غربال‌گری (\lr{Vetting})}\\
  \rl{آموزش حقوق بشر}\\[2pt]
  \rl{\scriptsize الگو: قانون اساسی آلمان}
};

\draw[maharrow=MahestanBlue] (center) -- (p1);
\draw[maharrow=MahestanRed] (center) -- (p2);
\draw[maharrow=MahestanOrange] (center) -- (p3);
\draw[maharrow=MahestanGreen] (center) -- (p4);
\draw[-Stealth, thick, MahestanBlue!50, dashed] (p1) -- (p2)
  node[midway, right]{\tiny\rl{مبنای محاکمات}};
\draw[-Stealth, thick, MahestanOrange!50, dashed] (p1) -- (p3)
  node[midway, above]{\tiny\rl{شناسایی قربانیان}};
\draw[-Stealth, thick, MahestanGreen!50, dashed] (p2) -- (p4)
  node[midway, right]{\tiny\rl{بازدارندگی}};
\end{tikzpicture}
}%
\end{center}
\vspace{0.5\baselineskip}

\section{مدل پنج‌ستونی جامع}

\begin{notebox}[title={\rl{مدل جامع (\lr{Holistic Model}) --- با ستون مستقل آشتی ملی}}]
مدل پنج‌ستونی با افزودن \textbf{«آشتی ملی و حافظه‌ی تاریخی»} به‌عنوان ستون مستقل، بر بُعد اجتماعی‌فرهنگی تأکید بیشتری دارد. این مدل در کلمبیا، سیرالئون و تیمور شرقی اجرا شده است.
\end{notebox}

\vspace{0.5\baselineskip}
\begin{center}
\resizebox{0.95\linewidth}{!}{%
\begin{tikzpicture}[
  every node/.style={align=right, font=\small},
  node distance=0.7cm and 1.5cm
]
\node[
  rectangle, rounded corners=8pt,
  draw=MahestanPurple, fill=MahestanLightPurple,
  text width=4.5cm, inner sep=10pt, font=\bfseries\small
] (center5) {
  \rl{عدالت انتقالی ایران}\\[3pt]
  \rl{\footnotesize مدل پنج‌ستونی جامع}\\
  \rl{\footnotesize (\lr{Holistic Model})}
};

\node[mahbox=MahestanBlue, above=2cm of center5, xshift=-4.5cm] (q1) {
  \textbf{\rl{ستون ۱: حقیقت‌یابی}}\\[2pt]
  \rl{کمیسیون حقیقت‌یابی ملی}\\
  \rl{مستندسازی نظام‌مند}\\
  \rl{جمع‌آوری شهادت‌ها}\\
  \rl{جلسات علنی}
};
\node[mahbox=MahestanRed, above=2cm of center5, xshift=-1.5cm] (q2) {
  \textbf{\rl{ستون ۲: عدالت کیفری}}\\[2pt]
  \rl{دادگاه‌های داخلی ویژه}\\
  \rl{همکاری با \lr{ICC}}\\
  \rl{تفکیک سطوح مسئولیت}\\
  \rl{امکان عفو مشروط}
};
\node[mahbox=MahestanOrange, above=2cm of center5, xshift=1.5cm] (q3) {
  \textbf{\rl{ستون ۳: جبران خسارت}}\\[2pt]
  \rl{الف) فردی: غرامت مالی}\\
  \rl{ب) جمعی: بازسازی مناطق}\\
  \rl{ج) نمادین: عذرخواهی}
};
\node[mahbox=MahestanGold, above=2cm of center5, xshift=4.5cm] (q4) {
  \textbf{\rl{ستون ۴: آشتی ملی}}\\[2pt]
  \rl{گفتگوهای ملی}\\
  \rl{موزه‌ها و یادبودها}\\
  \rl{اصلاح تاریخ‌نگاری}\\
  \rl{روز ملی یادبود}
};
\node[mahbox=MahestanGreen, below=2cm of center5] (q5) {
  \textbf{\rl{ستون ۵: تضمین عدم تکرار}}\\[2pt]
  \rl{اصلاح بخش امنیتی}\\
  \rl{غربال‌گری مقامات}\\
  \rl{آموزش حقوق بشر}\\
  \rl{قانون اساسی دموکراتیک}
};

\draw[maharrow=MahestanBlue] (center5) -- (q1);
\draw[maharrow=MahestanRed] (center5) -- (q2);
\draw[maharrow=MahestanOrange] (center5) -- (q3);
\draw[maharrow=MahestanGold] (center5) -- (q4);
\draw[maharrow=MahestanGreen] (center5) -- (q5);
\draw[-Stealth, thick, MahestanBlue!40, dashed] (q1) -- (q2);
\draw[-Stealth, thick, MahestanGold!60, dashed] (q4) -- (q5)
  node[midway, left]{\tiny\rl{فرهنگ‌سازی}};
\end{tikzpicture}
}%
\end{center}
\vspace{0.5\baselineskip}

\section{مقایسه‌ی دو مدل و کدام برای ایران}

\vspace{0.5\baselineskip}
\begin{adjustbox}{max width=\linewidth}
\begin{tabular}{@{}p{3.5cm}p{5.5cm}p{5.5cm}@{}}
\toprule
\textbf{معیار} & \textbf{مدل ۴ ستونی (سازمان ملل)} & \textbf{مدل ۵ ستونی (جامع)} \\
\midrule
\rowcolor{MahestanLightBlue!30}
\textbf{مرجع اصلی} &
  \lr{UN Guidance Note} (2010) &
  \lr{ICTJ}؛ \lr{Pablo de Greiff}؛ تجربه‌ی کلمبیا \\
\addlinespace
\textbf{تعداد ستون‌ها} & ۴ & ۵ \\
\addlinespace
\rowcolor{MahestanLightBlue!30}
\textbf{جایگاه آشتی} & زیرمجموعه‌ی سایر ستون‌ها & \textbf{ستون مستقل و برابر} \\
\addlinespace
\textbf{تأکید اصلی} & حقوقی‌قانونی & حقوقی‌قانونی + \textbf{اجتماعی‌فرهنگی} \\
\addlinespace
\rowcolor{MahestanLightBlue!30}
\textbf{مناسب برای} & جوامع با شکاف‌های کمتر & جوامع با شکاف‌های عمیق \\
\addlinespace
\textbf{مزیت اصلی} & سادگی؛ اجماع بین‌المللی & جامعیت؛ توجه به بُعد فرهنگی \\
\addlinespace
\rowcolor{MahestanLightBlue!30}
\textbf{ضعف اصلی} & ممکن است بُعد اجتماعی نادیده گرفته شود & پیچیدگی بیشتر؛ منابع و زمان بیشتر \\
\addlinespace
\rowcolor{NodeBg}
\textbf{پیشنهاد برای ایران} &
  \multicolumn{2}{p{11cm}}{\rl{مجلس مهستان بر اساس مشاوره‌ی عمومی و بررسی تجربه‌ی جهانی، یکی از این دو مدل (یا ترکیبی از آن‌ها) را انتخاب خواهد کرد.}} \\
\bottomrule
\end{tabular}
\end{adjustbox}

% ============================================================
%  APPENDIX: ELECTORAL SYSTEM
% ============================================================
\chapter{ضمیمه‌ی ک: سیستم انتخاباتی مجلس مهستان و آینده}
\label{app:electoral}

\section{مقایسه‌ی سیستم‌های انتخاباتی}

\vspace{0.5\baselineskip}
\noindent
\begin{adjustbox}{max width=\linewidth}
\begin{tabular}{@{}p{2.5cm}p{2.8cm}p{2.8cm}p{2.8cm}p{2.8cm}@{}}
\toprule
\textbf{معیار} & \textbf{\lr{FPTP} (اکثریتی)} & \textbf{\lr{List PR} (نسبی)} & \textbf{\lr{MMP} (مختلط)} & \textbf{\lr{STV}} \\
\midrule
نمایندگی منطقه‌ای & بالا & پایین & متوسط-بالا & متوسط \\
\addlinespace
\rowcolor{MahestanLightGray}
نمایندگی اقلیت‌ها & ضعیف & بالا & بالا & بالا \\
\addlinespace
ثبات حکومتی & بالا & متوسط & متوسط-بالا & متوسط \\
\addlinespace
\rowcolor{MahestanLightGray}
پیچیدگی اجرایی & ساده & ساده & متوسط & پیچیده \\
\addlinespace
مناسب برای گذار & متوسط & \textbf{بالا} & متوسط & پایین \\
\addlinespace
\rowcolor{NodeBg}
\textbf{پیشنهاد} & --- & \textbf{مجلس مهستان} & \textbf{قانون اساسی نهایی} & --- \\
\bottomrule
\end{tabular}
\end{adjustbox}

\section{توجیه سیستم \lr{MMP} برای آینده}

\begin{infobox}[title={\rl{چرا \lr{MMP} برای قانون اساسی نهایی؟}}]
سیستم \lr{MMP} (مختلط موازی‌جبرانی --- الگوی آلمانی با دو رأی) مزایای هر دو رویکرد
را ترکیب می‌کند:
\begin{itemize}
  \item \textbf{رأی اول}: رأی به نماینده‌ی حوزه‌ی محلی --- نمایندگی منطقه‌ای قوی
  \item \textbf{رأی دوم}: رأی به فهرست حزبی ملی --- جبران‌کننده و تضمین نسبیت
\end{itemize}
این سیستم در \textbf{آلمان} (بیش از ۷۰ سال)، \textbf{نیوزیلند}، \textbf{اسکاتلند}
و چندین کشور دیگر با موفقیت اجرا شده است.

\textbf{چرا نه برای مجلس مهستان؟} اجرای \lr{MMP} نیازمند زیرساخت فنی بیشتر
(ثبت احزاب، محاسبه‌ی جبرانی) و آموزش رأی‌دهندگان است که در هفته‌های اول گذار
فراهم نیست. سیستم فهرستی نسبی‌تر برای اولین انتخابات پس از گذار مناسب‌تر است.
\end{infobox}

% ============================================================
%  APPENDIX: INTERNATIONAL ENGAGEMENT
% ============================================================
\chapter{ضمیمه‌ی و: راهبرد تعامل بین‌المللی}
\label{app:international}

\section{نقش‌های بین‌المللی مورد نیاز}

\vspace{0.5\baselineskip}
\noindent
\begin{adjustbox}{max width=\linewidth}
\begin{tabular}{@{}p{4cm}p{9cm}@{}}
\toprule
\textbf{نهاد} & \textbf{نقش مورد انتظار} \\
\midrule
\rowcolor{MahestanLightGray}
سازمان ملل (\lr{UNDP}، \lr{OHCHR}) &
  نظارت بر انتخابات؛ کمک فنی در تدوین قانون اساسی؛ حمایت از عدالت انتقالی \\
اتحادیه اروپا &
  حمایت اقتصادی؛ مشاوره‌ی حقوقی و قضایی؛ مذاکرات تجاری \\
\rowcolor{MahestanLightGray}
صندوق بین‌المللی پول و بانک جهانی &
  بسته‌ی حمایتی اقتصادی اضطراری؛ مشاوره‌ی اصلاحات مالی و بانکی \\
کمیسیون ونیز &
  مشاوره‌ی تخصصی حقوقی در تدوین قانون اساسی \\
\rowcolor{MahestanLightGray}
\lr{IDEA} بین‌المللی &
  کمک فنی در طراحی سیستم انتخاباتی \\
دیوان بین‌المللی کیفری (\lr{ICC}) &
  همکاری در رسیدگی به جنایات علیه بشریت \\
\rowcolor{MahestanLightGray}
دیاسپورای ایرانی &
  مشارکت در رأی‌گیری؛ سرمایه‌گذاری؛ انتقال دانش و تخصص \\
\bottomrule
\end{tabular}
\end{adjustbox}

% ============================================================
%  APPENDIX: DUAL CONSTITUTION
% ============================================================
\chapter{ضمیمه‌ی ز: مدل دو قانون اساسی}
\label{app:dualconst}

\section{مقایسه‌ی قانون اساسی موقت و نهایی}

\vspace{0.5\baselineskip}
\noindent
\begin{adjustbox}{max width=\linewidth}
\begin{tabular}{@{}p{3.5cm}p{5cm}p{5cm}@{}}
\toprule
\textbf{معیار} & \textbf{قانون اساسی موقت} & \textbf{قانون اساسی نهایی} \\
\midrule
مرجع تدوین & کمیسیون ویژه‌ی مجلس مهستان & مجلس/کمیسیون مؤسسان منتخب \\
\addlinespace
\rowcolor{MahestanLightGray}
زمان تدوین & ۶ تا ۸ هفته & ۶ تا ۱۲ ماه \\
\addlinespace
عمق و تفصیل & حداقلی و عملیاتی (۴۰--۶۰ اصل) & جامع (۱۵۰--۲۵۰ اصل) \\
\addlinespace
\rowcolor{MahestanLightGray}
تصویب در مجلس & \lr{$\frac{2}{3}$} آرای مجلس مهستان & \lr{$\frac{2}{3}$} آرای مجلس مؤسسان \\
\addlinespace
رفراندوم & بله (۵۰\%+۱) & بله (۵۰\%+۱) \\
\addlinespace
\rowcolor{MahestanLightGray}
قابلیت تغییر & با رأی \lr{$\frac{2}{3}$} مجلس مهستان & فقط با رفراندوم یا فوق‌اکثریت \lr{$\frac{4}{5}$} \\
\addlinespace
مدت اعتبار & تا تصویب قانون اساسی نهایی (حداکثر ۲ سال) & دائمی \\
\addlinespace
\rowcolor{SuccessBg}
\textbf{اهمیت} & جلوگیری از خلأ حقوقی + سرعت & ساختار نهایی و دائمی کشور \\
\bottomrule
\end{tabular}
\end{adjustbox}

% ============================================================
%  APPENDIX: GLOSSARY
% ============================================================
\chapter{ضمیمه‌ی ر: واژه‌نامه‌ی تخصصی}
\label{app:glossary}

\begin{adjustbox}{max width=\linewidth}
\begin{longtable}{@{}p{4cm}p{10cm}@{}}
\toprule
\textbf{واژه} & \textbf{توضیح} \\
\midrule
\endhead
\bottomrule
\endfoot
\rl{گذار دموکراتیک} &
  فرایند تبدیل از یک نظام اقتدارگرا به یک نظام دموکراتیک. شامل مرحله‌ی
  «آزادسازی» (کاهش سرکوب)، «گذار» (تثبیت نهادهای دموکراتیک)
  و «تحکیم» (عمیق شدن دموکراسی). \\
\addlinespace
\rowcolor{MahestanLightGray}
\rl{مجلس مهستان} &
  نهاد اصلی گذار: مجلس منتخب مردم که مسئولیت هدایت فرایند گذار
  (قانون‌گذاری، نهادسازی، قانون اساسی) را بر عهده دارد. \\
\addlinespace
\rl{قانون اساسی‌گرایی دو مرحله‌ای} &
  مدلی که در آن یک قانون اساسی موقت (سریع، چارچوب‌ساز)
  و سپس یک قانون اساسی نهایی (جامع، دائمی) تدوین می‌شود.
  نمونه: آفریقای جنوبی (۱۹۹۴ و ۱۹۹۶). \\
\addlinespace
\rowcolor{MahestanLightGray}
\lr{SSR} &
  \lr{Security Sector Reform}: اصلاح بخش امنیتی؛ فرایند تحول
  نیروهای نظامی، پلیس و دستگاه اطلاعاتی از ابزار سرکوب به
  خدمتگزار قانون و مردم. \\
\addlinespace
\rl{عدالت انتقالی} &
  مجموعه‌ای از فرایندها و سازوکارها که جوامع پس از دوره‌ای از
  سرکوب یا درگیری برای پردازش گذشته به کار می‌برند؛ شامل
  حقیقت‌یابی، عدالت کیفری، جبران خسارت، و اصلاح نهادی. \\
\addlinespace
\rowcolor{MahestanLightGray}
\lr{MMP} &
  \lr{Mixed Member Proportional}: سیستم مختلط موازی‌جبرانی؛
  هر رأی‌دهنده دو رأی دارد: یکی برای نماینده‌ی محلی و یکی برای
  فهرست حزبی؛ کرسی‌های جبرانی تضمین می‌کنند نتایج نهایی نسبی باشد. \\
\addlinespace
\rl{بازوی تکنوکراتیک} &
  مرکز تحقیقات مجلس مهستان؛ مسئول تهیه‌ی گزارش‌های فنی‌تخصصی
  و ارائه‌ی گزینه‌های سیاستی بدون توصیه‌ی یکجانبه. \\
\addlinespace
\rowcolor{MahestanLightGray}
\rl{بازوی دموکراتیک} &
  واحد ارتباط مردمی و دموکراسی مستقیم؛ مسئول نظرسنجی، جلسات
  عمومی، مدیریت طومارها و پیوند مستقیم با مردم. \\
\addlinespace
\lr{Vetting} (غربال‌گری) &
  فرایند بررسی سوابق مقامات و نیروهای امنیتی برای شناسایی و
  کنار گذاشتن عاملان نقض حقوق بشر از مناصب عمومی. \\
\addlinespace
\rowcolor{MahestanLightGray}
\rl{دیاسپورا} &
  جمعیت ایرانیان مقیم خارج از کشور. این سند بر حق مشارکت
  دیاسپورا در انتخابات گذار تأکید دارد. \\
\end{longtable}
\end{adjustbox}

% ============================================================
%  CLOSING
% ============================================================
\chapter*{خاتمه}
\addcontentsline{toc}{chapter}{خاتمه}

\begin{summarybox}
این سند نه یک فرمان است و نه یک پیمان‌نامه‌ی قطعی. \textbf{دعوتی است برای گفتگو.}

طرح مجلس مهستان یک پیشنهاد است --- چارچوبی که با مشارکت همگانی بهتر خواهد شد.
هر نقد سازنده، هر پیشنهاد بهتر، و هر تجربه‌ی تاریخی که از قلم افتاده، این سند را
غنی‌تر می‌کند.

\begin{center}
{\large\color{MahestanDarkBlue}\bfseries
تصمیم نهایی درباره‌ی آینده‌ی ایران\\[0.3em]
فقط و فقط با مردم ایران است.}
\end{center}
\end{summarybox}

\vspace{2cm}
\begin{flushright}
\small\color{MahestanGray}
آخرین ویرایش: اسفند ۱۴۰۴ \lr{/} فوریه‌ی ۲۰۲۶ --- نسخه‌ی ۲.۰

\vspace{0.3cm}
جزئیات بیشتر در کتاب پیشنهادی «شکوفایی: ایران از بحران تا بالندگی» ارائه خواهد شد.
\end{flushright}

% ── Bibliography placeholder ────────────────────────────────
\backmatter
\chapter*{منابع و مراجع}
\addcontentsline{toc}{chapter}{منابع و مراجع}

\begin{itemize}[itemsep=4pt]
  \item \rl{لیپهارت، آرند.}
    \textit{\rl{الگوهای دموکراسی: اشکال حکومت و عملکرد در ۳۶ کشور.}}
    \lr{Lijphart, Arend (1999). \textit{Patterns of Democracy.} Yale University Press.}

  \item \rl{هانتینگتون، ساموئل.}
    \lr{Huntington, Samuel (1991). \textit{The Third Wave: Democratization in the Late 20th Century.}}

  \item \rl{اودانل و اشمیتر.}
    \lr{O'Donnell \& Schmitter (1986). \textit{Transitions from Authoritarian Rule.} Johns Hopkins.}

  \item \rl{پرز‌دیاز، ویکتور.}
    \lr{Pérez-Díaz, Víctor (1993). \textit{The Return of Civil Society: The Emergence of Democratic Spain.}}

  \item \rl{کمیسیون ونیز.}
    \lr{Venice Commission (2014). \textit{Rule of Law Checklist.} Council of Europe.}

  \item \rl{کمیسیون حقیقت و آشتی آفریقای جنوبی.}
    \lr{Truth and Reconciliation Commission of South Africa (2003). \textit{Final Report, Volumes 1--7.}}

  \item \rl{ایدئا بین‌المللی.}
    \lr{International IDEA (2021). \textit{Electoral System Design: The New International IDEA Handbook.}}

  \item \rl{سازمان ملل.}
    \lr{United Nations (2010). \textit{Guidance Note of the Secretary-General: United Nations Approach to Transitional Justice.}}
\end{itemize}